% Vietnamese translation for OI Public Library
% Source-File: IOI中国国家候选队论文/2007/day1/1.高逸涵《与圆有关的离散化方法》.ppt
\documentclass[11pt]{article}
\usepackage{oipl-vn}

\newcommand{\Oh}{\mathcal{O}}

\title{Bản trình chiếu: Phương pháp rời rạc hóa liên quan đến đường tròn}
\author{Gao Yihan\\Trường Phổ thông trực thuộc Đại học Thanh Hoa, Bắc Kinh}
\date{2007}

\begin{document}
\maketitle

\origtitle{Discretization methods related to circles}
\sourcepath{IOI China candidate team papers / 2007 / day 1 / Gao Yihan.ppt}

\section{Dẫn nhập}

Rời rạc hóa là một thủ pháp quen thuộc trong hình học tính toán. Với bài toán
kinh điển ``cho \(N\) hình chữ nhật, tính diện tích hợp'', ta lấy các biên
trên và dưới của hình chữ nhật làm đường chia. Sau khi chia, trong mỗi dải
ngang, phần bị phủ trên trục hoành có cấu trúc ổn định; diện tích của từng
phần có thể tính rồi cộng lại.

Vấn đề mới của bài trình chiếu là: nếu đối tượng không còn là hình có biên gấp
khúc, mà là các đường cong như đường tròn, thì rời rạc hóa có còn hữu dụng hay
không? Đường cong không có các ``điểm ngoặt'' tự nhiên như đa giác, nên cách
lấy tọa độ biên theo lối cũ dường như mất tác dụng.

Quan sát quan trọng là ta không nhất thiết đòi hỏi mỗi ô sau khi chia có tính
chất quá mạnh. Chỉ cần mỗi vùng có một thuộc tính thô đủ đơn giản để xử lý,
việc chia vẫn tiếp tục được. Với đường tròn, có thể thêm các đường cắt ngang
hoặc các đường cắt qua sự kiện hình học; vùng thu được thường gồm các mảnh
đơn giản như hình thang và cung tròn.

\section{Khuôn mẫu rời rạc hóa đường tròn}

Các slide nêu ba bước chung:
\begin{enumerate}
  \item Dựa vào bài toán, chia các đường tròn trên mặt phẳng thành một số vùng
  sao cho trong mỗi vùng, dữ liệu có thuộc tính thô tương đối đơn giản. Thông
  thường dùng các đường thẳng để chia.
  \item Dựa vào thuộc tính ấy, chọn thuật toán cụ thể trong từng vùng và tính
  kết quả cục bộ.
  \item Tổng hợp kết quả của mọi vùng để thu được đáp án cuối cùng.
\end{enumerate}

Điểm mấu chốt là mức độ chia. Nếu chia quá mịn, tính toán trong mỗi vùng rất
dễ nhưng số vùng và chi phí hợp nhất tăng mạnh. Nếu chia quá thô, cấu trúc bên
trong mỗi vùng vẫn phức tạp và không còn thuật toán đơn giản. Vì vậy rời rạc
hóa không chỉ là ``đổi tọa độ thực thành chỉ số mảng'', mà là chọn đúng các sự
kiện làm quan hệ hình học thay đổi.

\section{Ví dụ 1: \textit{Dolphin Pool}}

Bài \textit{Dolphin Pool} cho \(N\le 20\) đường tròn, biết tâm và bán kính.
Không có hai đường tròn tiếp xúc, và không có tâm đường tròn nào nằm trong một
đường tròn khác. Cần đếm số miền đóng nằm ngoài tất cả các hình tròn.

\subsection{Cách thử dựa trên flood fill}

Vì tọa độ nhỏ và input/output đều là số nguyên, cách nghĩ đầu tiên là dựng một
lưới điểm trên mặt phẳng, đánh dấu mọi điểm nằm trong hình tròn, rồi DFS trên
các điểm còn lại để đếm thành phần liên thông; số thành phần trừ \(1\) là số
miền đóng. Cách này dễ viết nhưng phụ thuộc hoàn toàn vào mật độ điểm lấy mẫu.
Muốn chính xác hơn phải tăng mật độ lưới, còn khi tăng mật độ thì thời gian và
bộ nhớ vượt quá giới hạn.

\subsection{Lần rời rạc hóa thứ nhất}

Xét từng hàng ngang của lưới. Một đường tròn cắt một hàng ngang thành một đoạn
liên tục, hoặc không cắt gì. Do đó thay vì ghi từng điểm trên hàng, ta chỉ ghi
các đoạn bị phủ bởi các đường tròn. Với đường ngang \(y=h\), nếu đường tròn có
tâm \((x_0,y_0)\) và bán kính \(r\), đoạn phủ là
\[
  \left[
    x_0-\sqrt{r^2-(y_0-h)^2},
    x_0+\sqrt{r^2-(y_0-h)^2}
  \right]
\]
khi \(|y_0-h|\le r\).

Sau đó sắp xếp và gộp các đoạn bị phủ. Bài toán gộp đoạn có thể làm tuyến tính
sau khi đã sắp xếp: mỗi lần đưa một đoạn mới vào, nếu nó giao đoạn trên đỉnh
ngăn xếp thì gộp lại, lặp cho đến khi không còn giao nữa. Phần bù của các đoạn
đã gộp là các khoảng chưa bị phủ trên hàng ngang.

Mỗi khoảng chưa bị phủ được xem như một đỉnh của đồ thị. Hai đỉnh ở hai hàng
kề nhau được nối cạnh nếu hai khoảng tương ứng giao nhau. Đếm số thành phần
liên thông của đồ thị này sẽ cho số miền. Để giảm bộ nhớ, không cần giữ toàn
bộ đồ thị; có thể quét theo từng lớp hàng ngang và dùng DSU để nối các khoảng
của hai lớp liên tiếp. Một phép gộp thành công làm số thành phần giảm đi một.

Cách này cải thiện đáng kể tốc độ và độ chính xác, nhưng vẫn có nhược điểm:
các hàng ngang hầu hết chỉ lặp lại cùng quan hệ kế thừa giữa các đoạn. Với dữ
liệu cực hạn, độ chính xác vẫn chưa đủ ổn định.

\subsection{Lần rời rạc hóa thứ hai}

Thay vì duyệt mọi hàng, ta chỉ xét các thời điểm quan hệ giữa các khoảng có
thể thay đổi. Các slide gọi đó là các \emph{sự kiện điểm}. Những sự kiện quan
trọng gồm:
\begin{itemize}
  \item một đường tròn mới xuất hiện, làm một khoảng bị tách làm hai;
  \item đi qua đáy của một đường tròn, hai khoảng hợp lại;
  \item hai đường tròn giao nhau, một khoảng co dần rồi biến mất;
  \item hai đường tròn giao nhau, một khoảng mới được sinh ra.
\end{itemize}

Nếu mô phỏng trực tiếp sự sinh-trưởng-mất của mọi khoảng thì cài đặt khá phức
tạp. Cách ``lười'' được slide đề xuất là: trước và sau mỗi sự kiện lấy một
đường ngang đại diện, xử lý hai đường này bằng phương pháp gộp đoạn của lần
rời rạc hóa thứ nhất; giữa hai sự kiện, quan hệ kế thừa của các khoảng được
xem là một-một. Nhờ vậy ta vẫn bắt được mọi thay đổi quan trọng mà không cần
theo dõi từng hàng. Độ phức tạp cuối cùng của bài này được nêu là
\(\Oh(N^3)\).

\section{Ví dụ 2: \textit{Empire Strikes Back}}

Bài \textit{Empire Strikes Back} có một đường tròn lớn tâm \((0,0)\) và
\(N\le 300\) đường tròn nhỏ có tâm nằm trong đường tròn lớn. Mọi đường tròn
nhỏ có cùng bán kính. Cần tìm bán kính nhỏ nhất để các đường tròn nhỏ phủ kín
đường tròn lớn.

Vì đáp án là một số thực, ta nhị phân bán kính. Với một bán kính đã chọn, việc
còn lại là kiểm tra đường tròn lớn có được phủ kín hay không.

Nếu áp dụng máy móc cách rời rạc hóa của ví dụ 1, độ phức tạp là
\(\Oh(N^3k)\), trong đó \(k\) là số vòng nhị phân; cách này quá chậm. Nguyên
nhân là sau khi nhị phân, bài toán đã trở thành bài toán quyết định ``có phủ
kín hay không'', nhưng ta vẫn dùng tư duy đếm miền của ví dụ trước.

Slide chuyển sang quan sát các điểm đặc biệt trên một đường ngang: các đầu mút
của những đoạn được phủ và các điểm rất gần hai phía của chúng. Nếu mọi điểm
đặc biệt đều bị phủ thì cả đường ngang được phủ. Nhìn toàn cục, các điểm đặc
biệt này tương ứng với việc dịch mỗi đường tròn sang trái và sang phải một
khoảng rất nhỏ. Phần mới sinh ra bởi phép dịch có thể dùng để nhận diện biên
chưa phủ.

Từ đó bài toán được chuyển thành: với mỗi đường tròn nhỏ, xét phần biên của
đường tròn lớn nằm trên nó; kiểm tra xem phần biên ấy có được các đường tròn
nhỏ còn lại phủ hết hay không. Đây lại là bài toán gộp các khoảng trên một
đối tượng một chiều, tương tự lần rời rạc hóa đầu tiên trong ví dụ 1.

Độ phức tạp được trình bày là
\[
  \Oh(N^2\log N\cdot k),
\]
gồm nhị phân \(k\), sắp xếp các khoảng \(\Oh(N\log N)\), gộp đoạn tuyến tính,
và lặp qua \(N\) đường tròn. Một vài tối ưu có thể đưa biểu thức về
\(\Oh(N^2\log N+N\log N\cdot k)\) theo phân tích của slide.

\section{Tổng kết}

Rời rạc hóa trong các bài toán có đường tròn không chỉ là chia mặt phẳng theo
tọa độ ngang dọc. Cốt lõi là biến ``nhiều điểm liên tục'' thành ``một số sự
kiện hữu hạn'' nhưng vẫn giữ đúng quan hệ phủ, giao và liên thông cần dùng.
Với cùng một bài toán, hai cách rời rạc hóa khác nhau có thể dẫn tới độ phức
tạp rất khác nhau. Khi chọn đúng mức chia, ta vừa giảm độ phức tạp cài đặt vừa
đạt được các chặn như \(\Oh(N^3)\) hoặc \(\Oh(N^2\log N\cdot k)\).

\end{document}
